% 112-45(인쇄 112쪽) — 개념 34 45번의 표준화 그림: 왼쪽은 정규분포 N(40, 3^2) 곡선에서 37~49 를 색칠한 것(대칭축 40 은 점선),
% 오른쪽은 같은 부분을 표준정규분포 곡선 위에 옮겨 a, b 로 표시한 것. 스캔 화질이 낮아 TikZ 로 다시 그림.
% 라벨은 원문에 있는 것만(37, 40, 49, x / a, 0, b, z)이고 a, b 의 값은 그림에 적지 않는다.
\begin{tikzpicture}[x=0.56cm, y=0.95cm, line width=0.55pt, font=\small]
  \fill[magenta!18] plot[domain=-1:3, samples=80, smooth] ({\x},{exp(-\x*\x/2)}) -- (3,0) -- (-1,0) -- cycle;
  \draw[->, >=stealth, line width=0.7pt] (-3.6,0) -- (3.9,0) node[below=1pt] {$x$};
  \draw plot[domain=-3.2:3.2, samples=110, smooth] ({\x},{exp(-\x*\x/2)});
  \draw[line width=0.45pt] (-1,0) -- (-1,{exp(-0.5)});
  \draw[line width=0.45pt] (3,0) -- (3,{exp(-4.5)});
  \draw[densely dotted, line width=0.4pt] (0,0) -- (0,1);
  \node[below=1pt, font=\footnotesize] at (-1,0) {$37$};
  \node[below=1pt, font=\footnotesize] at (0,0) {$40$};
  \node[below=1pt, font=\footnotesize] at (3,0) {$49$};
  \node[below=1pt, align=center] at (0.2,-0.55) {정규분포 $\mathrm{N}(40,\,3^2)$을\\ 따르는 확률변수 $X$};
  \draw[->, >=stealth, line width=1.2pt, blue!45] (4.9,0.7) -- (6.6,0.7);
  \begin{scope}[shift={(11.2,0)}]
    \fill[magenta!18] plot[domain=-1:3, samples=80, smooth] ({\x},{exp(-\x*\x/2)}) -- (3,0) -- (-1,0) -- cycle;
    \draw[->, >=stealth, line width=0.7pt] (-3.6,0) -- (3.9,0) node[below=1pt] {$z$};
    \draw plot[domain=-3.2:3.2, samples=110, smooth] ({\x},{exp(-\x*\x/2)});
    \draw[line width=0.45pt] (-1,0) -- (-1,{exp(-0.5)});
    \draw[line width=0.45pt] (3,0) -- (3,{exp(-4.5)});
    \draw[densely dotted, line width=0.4pt] (0,0) -- (0,1);
    \node[below=1pt, font=\footnotesize] at (-1,0) {$a$};
    \node[below=1pt, font=\footnotesize] at (0,0) {$0$};
    \node[below=1pt, font=\footnotesize] at (3,0) {$b$};
    \node[below=1pt, align=center] at (0.2,-0.55) {표준정규분포 $\mathrm{N}(0,\,1)$을\\ 따르는 확률변수 $Z$};
  \end{scope}
\end{tikzpicture}
